\documentclass[11pt, a4paper]{article}

% ── Packages ─────────────────────────────────────────────────────────────────
\usepackage[T1]{fontenc}
\usepackage[utf8]{inputenc}
\usepackage[french]{babel}
\usepackage[margin=2.2cm]{geometry}
\usepackage{amsmath, amssymb}
\usepackage{booktabs}
\usepackage{enumitem}
\usepackage{titlesec}
\usepackage[expansion=false]{microtype}
\usepackage[table]{xcolor}
\usepackage{tikz}
\usetikzlibrary{shapes, arrows.meta, positioning, calc}
\usepackage{fancybox}
\usepackage{tcolorbox}
\tcbuselibrary{skins, breakable}
\usepackage{multicol}
\usepackage{multirow}
\usepackage{float}

% ── Couleurs ─────────────────────────────────────────────────────────────────
\definecolor{accent}{HTML}{2196F3}
\definecolor{fair}{HTML}{4CAF50}
\definecolor{alertc}{HTML}{d9534f}
\definecolor{neutral}{HTML}{FF9800}
\definecolor{light}{HTML}{F5F5F5}
\definecolor{darkbg}{HTML}{263238}

% ── Boxes pédagogiques ──────────────────────────────────────────────────────
\newtcolorbox{questionbox}[1][]{%
    colback=accent!8, colframe=accent!60, fonttitle=\bfseries,
    title={#1}, arc=2mm, boxrule=0.5pt, left=4pt, right=4pt,
    top=3pt, bottom=3pt, breakable}

\newtcolorbox{leconbox}[1][]{%
    colback=fair!8, colframe=fair!60, fonttitle=\bfseries,
    title={#1}, arc=2mm, boxrule=0.5pt, left=4pt, right=4pt,
    top=3pt, bottom=3pt, breakable}

\newtcolorbox{attentionbox}[1][]{%
    colback=alertc!8, colframe=alertc!60, fonttitle=\bfseries,
    title={#1}, arc=2mm, boxrule=0.5pt, left=4pt, right=4pt,
    top=3pt, bottom=3pt, breakable}

\newtcolorbox{codebox}{%
    colback=light, colframe=darkbg!40, arc=1mm, boxrule=0.4pt,
    left=4pt, right=4pt, top=2pt, bottom=2pt,
    fontupper=\small\ttfamily}

% ── Mise en page ─────────────────────────────────────────────────────────────
\setlength{\parskip}{0.4em}
\setlength{\parindent}{0pt}
\titlespacing*{\section}{0pt}{1.2em}{0.5em}
\titlespacing*{\subsection}{0pt}{0.8em}{0.3em}
\setlist{nosep, leftmargin=1.5em}

% ══════════════════════════════════════════════════════════════════════════════
\title{\textbf{Parcours p\'{e}dagogique : IA Responsable\\[0.2em]
\large Du jeu de donn\'{e}es aux conclusions ---
explication pas \`{a} pas du notebook}}
\author{IA708 --- T\'{e}l\'{e}com Paris, 2026}
\date{}

\begin{document}
\maketitle
\vspace{-1em}

\begin{abstract}
\noindent Ce document accompagne le notebook \texttt{responsiveAI-german-credit.ipynb} et d\'{e}taille chaque \'{e}tape du parcours d'analyse. L'objectif : \'{e}valuer un mod\`{e}le de scoring cr\'{e}dit sous trois angles --- \textbf{\'{e}quit\'{e}}, \textbf{interpr\'{e}tabilit\'{e}} et \textbf{quantification d'incertitude} (qui remplace la robustesse adversariale, sur conseil de l'enseignante) --- en partant des donn\'{e}es brutes et en progressant vers des conclusions argument\'{e}es. Stack utilis\'{e} : \texttt{scikit-learn}, \texttt{fairlearn}, \texttt{shap} ; reproductibilit\'{e} via \texttt{papermill}. Ce parcours suit la structure du cours IA708.
\end{abstract}

\tableofcontents
\vspace{1em}

% ══════════════════════════════════════════════════════════════════════════════
\section{Vue d'ensemble du parcours}
\label{sec:overview}

Avant d'entrer dans le code, comprenons la \textbf{logique du parcours}. Chaque \'{e}tape r\'{e}pond \`{a} une question qui motive naturellement la suivante :

\begin{center}
\begin{tikzpicture}[
    node distance=0.4cm and 0.9cm,
    block/.style={rectangle, draw, rounded corners=3pt, minimum height=0.85cm,
                  minimum width=2.4cm, text centered, font=\small,
                  text width=2.3cm, align=center},
    question/.style={font=\scriptsize\itshape, text=accent!80!black, text width=3cm, align=center},
    arrow/.style={-Stealth, thick, accent!70}
]
\node[block, fill=accent!12] (data) {1.~Donn\'{e}es\newline \scriptsize 1000 lignes};
\node[block, fill=accent!12, right=of data] (explore) {2.~Exploration\newline \scriptsize biais?};
\node[block, fill=accent!12, right=of explore] (prep) {3.~Pr\'{e}paration\newline \scriptsize split, encoding};
\node[block, fill=accent!12, right=of prep] (model) {4.~Mod\`{e}le\newline \scriptsize r\'{e}g. logist.};

\node[block, fill=fair!12, below=1.1cm of data] (metrics) {5.~M\'{e}triques\newline \scriptsize perf + \'{e}quit\'{e}};
\node[block, fill=fair!12, right=of metrics] (rw) {6.~Reweighing\newline \scriptsize pr\'{e}-trait.};
\node[block, fill=fair!12, right=of rw] (pp) {7.~Seuils/grp\newline \scriptsize post-trait.};
\node[block, fill=fair!12, right=of pp] (comp) {8.~Comparer\newline \scriptsize 4 configs};

\node[block, fill=neutral!12, below=1.1cm of metrics] (shap) {9.~SHAP\newline \scriptsize interpr\'{e}t.};
\node[block, fill=neutral!12, right=of shap] (rob) {10.~Bootstrap\newline \scriptsize IC \& incertitude};
\node[block, fill=alertc!12, right=of rob] (ccl) {11.~Conclusion\newline \scriptsize le\c{c}ons};

\draw[arrow] (data) -- (explore);
\draw[arrow] (explore) -- (prep);
\draw[arrow] (prep) -- (model);
\draw[arrow] (model.south) -- ++(0,-0.35) -| (metrics.north);
\draw[arrow] (metrics) -- (rw);
\draw[arrow] (rw) -- (pp);
\draw[arrow] (pp) -- (comp);
\draw[arrow] (comp.south) -- ++(0,-0.35) -| (shap.north);
\draw[arrow] (shap) -- (rob);
\draw[arrow] (rob) -- (ccl);

% Questions entre les blocs
\node[question, above=0.15cm of explore] {Y a-t-il des biais ?};
\node[question, above=0.15cm of model] {Quelles performances ?};
\node[question, below=0.15cm of rw] {Peut-on corriger ?};
\node[question, below=0.15cm of comp] {Quel compromis ?};
\node[question, below=0.15cm of shap] {Pourquoi ces d\'{e}cisions ?};
\node[question, below=0.15cm of ccl] {Qu'a-t-on appris ?};
\end{tikzpicture}
\end{center}

Chaque fl\`{e}che correspond \`{a} une \textbf{transition motiv\'{e}e} : on ne passe \`{a} l'\'{e}tape suivante que parce que la pr\'{e}c\'{e}dente a soulev\'{e} une question. C'est cette logique narrative que nous allons d\'{e}rouler.

% ══════════════════════════════════════════════════════════════════════════════
\section{Les donn\'{e}es : point de d\'{e}part}
\label{sec:data}

\begin{questionbox}[\'{E}tape 1 --- Pourquoi ces donn\'{e}es ?]
Le \textbf{German Credit Dataset} (UCI Statlog, Hofmann 1994) est un cas d'\'{e}cole en IA responsable : 1\,000 demandes de cr\'{e}dit d\'{e}crites par 20 attributs. La cible est binaire : \textbf{d\'{e}faut de paiement} (1) ou non (0). Le scoring cr\'{e}dit est une d\'{e}cision \`{a} fort impact humain --- un refus injuste peut affecter la vie enti\`{e}re d'un demandeur.
\end{questionbox}

\paragraph{Chargement.} Le fichier \texttt{german.data} est un CSV s\'{e}par\'{e} par espaces, sans ent\^{e}te. La cible originale est cod\'{e}e 1 (bon) / 2 (mauvais) ; on la transforme en $\texttt{default} \in \{0, 1\}$.

\paragraph{Premiers constats.}

\begin{table}[H]
\centering
\begin{tabular}{lcc}
\toprule
\textbf{Caract\'{e}ristique} & \textbf{Valeur} & \textbf{Proportion} \\
\midrule
Bon payeur (d\'{e}faut = 0) & 700 & 70\,\% \\
Mauvais payeur (d\'{e}faut = 1) & 300 & 30\,\% \\
\midrule
Features num\'{e}riques & 7 & --- \\
Features cat\'{e}gorielles & 13 & --- \\
\bottomrule
\end{tabular}
\caption{Vue d'ensemble du dataset.}
\end{table}

Le d\'{e}s\'{e}quilibre 70/30 est important : un mod\`{e}le qui pr\'{e}dirait toujours \og bon \fg{} aurait 70\,\% d'accuracy. C'est pourquoi on utilisera des m\'{e}triques adapt\'{e}es (\S\ref{sec:metrics}).

% ══════════════════════════════════════════════════════════════════════════════
\section{Exploration : les biais pr\'{e}existent dans les donn\'{e}es}
\label{sec:explore}

\begin{questionbox}[\'{E}tape 2 --- Y a-t-il d\'{e}j\`{a} des biais ?]
Avant m\^{e}me de construire un mod\`{e}le, on examine si les donn\'{e}es refl\`{e}tent des disparit\'{e}s entre groupes. Si oui, tout mod\`{e}le performant les \textbf{reproduira}.
\end{questionbox}

\subsection{Extraction des attributs sensibles}

On d\'{e}finit deux attributs sensibles, conform\'{e}ment au cours \emph{fairness} :

\begin{itemize}
    \item \textbf{Genre} : extrait de la colonne \texttt{personal\_status\_sex}. Les codes A91, A93, A94 correspondent aux hommes ; A92, A95 aux femmes.
    \item \textbf{\^{A}ge} : binaris\'{e} au seuil de 25 ans. Adulte ($> 25$) vs jeune ($\leq 25$).
\end{itemize}

\subsection{R\'{e}partition des groupes}

\begin{table}[H]
\centering
\begin{tabular}{lccc}
\toprule
\textbf{Groupe} & \textbf{Effectif} & \textbf{Taux bon} & \textbf{Taux d\'{e}faut} \\
\midrule
\rowcolor{accent!6} Homme & 690 (69\,\%) & 72\,\% & 28\,\% \\
\rowcolor{accent!6} Femme & 310 (31\,\%) & 65\,\% & 35\,\% \\
\midrule
\rowcolor{fair!6} Adulte ($>25$) & 810 (81\,\%) & 72\,\% & 28\,\% \\
\rowcolor{fair!6} Jeune ($\leq 25$) & 190 (19\,\%) & 60\,\% & 40\,\% \\
\bottomrule
\end{tabular}
\caption{Taux de d\'{e}faut par groupe --- biais pr\'{e}existants dans les donn\'{e}es.}
\label{tab:bias}
\end{table}

\begin{leconbox}[Le\c{c}on cl\'{e}]
Les femmes ont un taux de d\'{e}faut plus \'{e}lev\'{e} que les hommes (+7 points), et les jeunes plus que les adultes (+12 points). Ces \'{e}carts ne sont pas n\'{e}cessairement \og vrais \fg{} --- ils peuvent refl\'{e}ter des biais historiques dans l'attribution des cr\'{e}dits. Un mod\`{e}le entra\^{i}n\'{e} sur ces donn\'{e}es reproduira et potentiellement amplifiera ces disparit\'{e}s.
\end{leconbox}

\subsection{Ce que le notebook montre}

Le notebook produit trois graphiques c\^{o}te \`{a} c\^{o}te :
\begin{enumerate}
    \item \textbf{Distribution de la cible} --- barre verte (bon) vs rouge (mauvais) : on voit le d\'{e}s\'{e}quilibre 70/30.
    \item \textbf{Distribution par genre} --- les hommes sont surrepr\'{e}sent\'{e}s (69\,\%).
    \item \textbf{Histogramme d'\^{a}ge} avec le seuil 25 ans --- la majorit\'{e} sont des adultes, les jeunes sont une minorit\'{e} (19\,\%).
\end{enumerate}

Ces visualisations motivent la suite : comment construire un mod\`{e}le qui soit performant \emph{et} \'{e}quitable ?

% ══════════════════════════════════════════════════════════════════════════════
\section{Pr\'{e}paration des donn\'{e}es}
\label{sec:prep}

\begin{questionbox}[\'{E}tape 3 --- Comment pr\'{e}parer les donn\'{e}es pour le mod\`{e}le ?]
Trois d\'{e}cisions cl\'{e}s : (1) retirer les attributs sensibles des features, (2) encoder les variables, (3) s\'{e}parer train/validation/test.
\end{questionbox}

\subsection{Retrait des attributs sensibles}

On retire \texttt{personal\_status\_sex} et \texttt{age\_in\_years} des features d'entra\^{i}nement.

\begin{attentionbox}[Attention : proxy]
Retirer la colonne ne suffit pas toujours. D'autres features peuvent \^{e}tre \textbf{corr\'{e}l\'{e}es} \`{a} l'attribut sensible et servir de \textbf{proxy} --- le mod\`{e}le peut alors discriminer \emph{indirectement}. On v\'{e}rifiera cela \`{a} l'\'{e}tape~9 (SHAP).
\end{attentionbox}

\subsection{Pr\'{e}traitement : sklearn ColumnTransformer}

\begin{itemize}
    \item \textbf{Num\'{e}riques} (6 colonnes) : \texttt{StandardScaler} (z-score), $\mu$ et $\sigma$ calcul\'{e}s sur le train uniquement.
    \item \textbf{Cat\'{e}gorielles} (12 colonnes) : \texttt{OneHotEncoder(handle\_unknown="ignore")}.
\end{itemize}

\noindent R\'{e}sultat : 56 features num\'{e}riques apr\`{e}s encodage.

Le pr\'{e}processeur est \textbf{fit\'{e} sur le train} et \textbf{appliqu\'{e}} au val/test --- pas de fuite d'information.

\begin{codebox}
preprocessor = ColumnTransformer([\\
\quad ("num", StandardScaler(), NUMERIC),\\
\quad ("cat", OneHotEncoder(handle\_unknown="ignore", sparse\_output=False), CATEG),\\
])\\
X\_tr = preprocessor.fit\_transform(features.iloc[idx\_tr])\\
X\_va = preprocessor.transform(features.iloc[idx\_va])\\
X\_te = preprocessor.transform(features.iloc[idx\_te])
\end{codebox}

\subsection{Split stratifi\'{e} 60/20/20}

La r\'{e}partition est stratifi\'{e}e par la cible : chaque sous-ensemble conserve le ratio 70/30. Le seed est fix\'{e} (\texttt{rng = default\_rng(42)}) pour la reproductibilit\'{e}.

\begin{table}[H]
\centering
\begin{tabular}{lccc}
\toprule
& \textbf{Train} & \textbf{Validation} & \textbf{Test} \\
\midrule
Taille & $\sim$600 & $\sim$200 & $\sim$200 \\
R\^{o}le & Entra\^{i}nement & Choix du seuil, PP & \'{E}valuation finale \\
\bottomrule
\end{tabular}
\end{table}

% ══════════════════════════════════════════════════════════════════════════════
\section{Mod\`{e}le de base : r\'{e}gression logistique sklearn}
\label{sec:model}

\begin{questionbox}[\'{E}tape 4 --- Quel mod\`{e}le, et comment l'entra\^{i}ner ?]
On utilise une r\'{e}gression logistique de \texttt{scikit-learn}. Pourquoi ce mod\`{e}le ? (1) C'est interpr\'{e}table (un coefficient par feature, formule SHAP exacte), (2) bien adapt\'{e} aux donn\'{e}es tabulaires, (3) calibr\'{e} par d\'{e}faut.
\end{questionbox}

\subsection{Le mod\`{e}le}

La r\'{e}gression logistique mod\'{e}lise la probabilit\'{e} de d\'{e}faut :

$$P(\text{d\'{e}faut}=1 \mid x) = \sigma(w^\top x + b) \quad \text{avec} \quad \sigma(z) = \frac{1}{1 + e^{-z}}$$

\noindent o\`{u} $w \in \mathbb{R}^d$ est le vecteur de poids et $b$ le biais.

\subsection{Entra\^{i}nement (sklearn)}

La perte est la \textbf{cross-entropy binaire} avec p\'{e}nalisation L2 :
$$\mathcal{L} = -\frac{1}{n}\sum_{i=1}^{n} s_i \Big[ y_i \ln p_i + (1-y_i) \ln(1 - p_i) \Big] + \frac{1}{2C} \|w\|^2$$

\noindent o\`{u} $s_i$ est le poids de l'exemple (utilis\'{e} pour le reweighing, $s_i = 1$ pour le baseline) et $C = 1$ est l'inverse de la r\'{e}gularisation (forme sklearn).

\begin{table}[H]
\centering
\begin{tabular}{ll}
\toprule
\textbf{Hyperparamètre} & \textbf{Valeur} \\
\midrule
Solver & \texttt{liblinear} (descente coordonn\'{e}e) \\
P\'{e}nalisation & L2, $C = 1$ \\
Max iterations & 1000 (convergence \`{a} $\sim$10 it\'{e}rations) \\
Random state & 42 \\
\bottomrule
\end{tabular}
\caption{Hyperparamètres de \texttt{LogisticRegression} sklearn.}
\end{table}

\subsection{V\'{e}rification de convergence (feedback \'{e}quipe)}

L'\'{e}quipe attendait une convergence \`{a} partir de 100 it\'{e}rations et une log-loss autour de 0.5. \textbf{Mesure :}

\begin{itemize}
    \item Log-loss train final : \textbf{0.43}
    \item Log-loss val : \textbf{0.56} (proche de 0.5 attendu sur d\'{e}s\'{e}quilibre 70/30)
    \item Convergence atteinte d\`{e}s 10 it\'{e}rations (\texttt{liblinear} converge tr\`{e}s vite sur ce dataset)
\end{itemize}

\subsection{Seuil de d\'{e}cision}

On ne fixe pas le seuil \`{a} 0.5. On balaye $\tau \in [0.1,\; 0.9]$ et on choisit celui qui maximise la \textbf{balanced accuracy} sur la validation :

$$\text{BalAcc} = \frac{1}{2}\left(\frac{TP}{TP+FN} + \frac{TN}{TN+FP}\right)$$

Ce choix est motiv\'{e} par le d\'{e}s\'{e}quilibre 70/30 : un seuil fixe \`{a} 0.5 favorise la classe majoritaire. R\'{e}sultat : $\tau = 0.18$ sur la validation.

% ══════════════════════════════════════════════════════════════════════════════
\section{M\'{e}triques : mesurer ce qui compte}
\label{sec:metrics}

\begin{questionbox}[\'{E}tape 5 --- Comment \'{e}valuer performance ET \'{e}quit\'{e} ?]
La performance seule ne suffit pas. Il faut aussi mesurer si le mod\`{e}le traite les groupes de mani\`{e}re \'{e}quitable. On d\'{e}finit les deux familles de m\'{e}triques utilis\'{e}es dans tout le reste du parcours.
\end{questionbox}

\subsection{M\'{e}triques de performance}

\begin{table}[H]
\centering
\begin{tabular}{lp{9cm}}
\toprule
\textbf{M\'{e}trique} & \textbf{Pourquoi ?} \\
\midrule
\textbf{AUC ROC} & Invariante au seuil. Mesure la capacit\'{e} \`{a} s\'{e}parer les classes. Ici calcul\'{e}e par la formule de Wilcoxon-Mann-Whitney. \\
\textbf{Balanced Acc.} & $\frac{1}{2}(\text{TPR}+\text{TNR})$. Adapt\'{e}e au d\'{e}s\'{e}quilibre (vs accuracy na\"{i}ve). \\
\textbf{ECE} & Expected Calibration Error. V\'{e}rifie que la confiance du mod\`{e}le correspond \`{a} la r\'{e}alit\'{e} (cf.~cours incertitude). \\
\bottomrule
\end{tabular}
\end{table}

\paragraph{ECE --- Expected Calibration Error.}
On d\'{e}coupe l'intervalle $[0,1]$ en $B$ bins par confiance, puis :
$$\text{ECE} = \sum_{b=1}^{B} \frac{|B_b|}{n}\, \big|\text{acc}(B_b) - \text{conf}(B_b)\big|$$
Un mod\`{e}le bien calibr\'{e} a ECE $\approx 0$. La r\'{e}gression logistique est naturellement bien calibr\'{e}e.

\subsection{M\'{e}triques d'\'{e}quit\'{e} (cf. cours fairness)}

Soit $S$ l'attribut sensible, $Y$ la cible, $\hat{Y}$ la pr\'{e}diction.

\begin{table}[H]
\centering
\begin{tabular}{lll}
\toprule
\textbf{Crit\`{e}re} & \textbf{D\'{e}finition} & \textbf{Signification} \\
\midrule
\textbf{Demographic Parity} & $\big|P(\hat{Y}{=}1 \mid S{=}0) - P(\hat{Y}{=}1 \mid S{=}1)\big|$ & M\^{e}me taux de s\'{e}lection \\
\textbf{Equal Opportunity} & $\big|\text{TPR}_{S=0} - \text{TPR}_{S=1}\big|$ & M\^{e}me rappel sur les vrais positifs \\
\bottomrule
\end{tabular}
\end{table}

Un mod\`{e}le parfaitement \'{e}quitable aurait $|\Delta\text{DP}| = 0$ et $|\Delta\text{EO}| = 0$. En pratique, c'est impossible simultan\'{e}ment (\S\ref{sec:impossibility}).

\subsection{\'{E}valuation du baseline}

Le notebook affiche les r\'{e}sultats du baseline :

\begin{table}[H]
\centering
\begin{tabular}{lcccc}
\toprule
\textbf{Attribut} & \textbf{Groupe} & \textbf{Taux s\'{e}lection} & \textbf{TPR} & \textbf{Effectif test} \\
\midrule
\multirow{2}{*}{Genre} & Homme & $\approx 0.41$ & $\approx 0.70$ & $\sim$140 \\
& Femme & $\approx 0.50$ & $\approx 0.74$ & $\sim$60 \\
\midrule
\multirow{2}{*}{\^{A}ge} & Adulte & $\approx 0.42$ & $\approx 0.77$ & $\sim$160 \\
& Jeune & $\approx 0.51$ & $\approx 0.59$ & $\sim$40 \\
\bottomrule
\end{tabular}
\caption{M\'{e}triques d'\'{e}quit\'{e} du baseline (valeurs approch\'{e}es, d\'{e}pendent du seed).}
\end{table}

\begin{attentionbox}[Constat]
Le baseline n'est pas \'{e}quitable. Pour l'\^{a}ge, le $|\Delta\text{EO}| \approx 0.18$ : les jeunes d\'{e}faillants sont beaucoup moins bien d\'{e}tect\'{e}s que les adultes. C'est \textbf{pr\'{e}occupant} --- cela justifie l'\'{e}tape suivante.
\end{attentionbox}

% ══════════════════════════════════════════════════════════════════════════════
\section{Reweighing : correction en pr\'{e}-traitement}
\label{sec:reweighing}

\begin{questionbox}[\'{E}tape 6 --- Peut-on corriger les biais sans changer le mod\`{e}le ?]
Le \textbf{reweighing} (cf.~cours fairness-mitigation) est une m\'{e}thode de pr\'{e}-traitement : on repondère les exemples d'entra\^{i}nement pour \'{e}quilibrer la relation entre $S$ et $Y$.
\end{questionbox}

\subsection{Principe}

L'id\'{e}e est simple : si un sous-groupe est sous-repr\'{e}sent\'{e} dans une classe, on augmente son poids. Formellement :

$$w_i = \frac{P(S = s_i) \cdot P(Y = y_i)}{P(S = s_i,\; Y = y_i)}$$

\noindent Ce poids rend $S \perp Y$ dans la distribution pond\'{e}r\'{e}e.

\subsection{Calcul concret sur un mini-exemple}

Imaginons un mini-train de 100 personnes, avec les comptages crois\'{e}s suivants :

\begin{table}[H]
\centering
\begin{tabular}{lccc}
\toprule
& $Y = 0$ (bon) & $Y = 1$ (d\'{e}faut) & \textbf{Total $S$} \\
\midrule
$S = $ Homme & 50 & 20 & 70 \\
$S = $ Femme & 20 & 10 & 30 \\
\midrule
\textbf{Total $Y$} & 70 & 30 & 100 \\
\bottomrule
\end{tabular}
\caption{Mini-exemple : effectifs crois\'{e}s.}
\end{table}

\noindent\textbf{Calcul des probabilit\'{e}s} :
$P(S{=}H) = 0{,}7$, $P(S{=}F) = 0{,}3$, $P(Y{=}0) = 0{,}7$, $P(Y{=}1) = 0{,}3$.

\noindent\textbf{Calcul des poids} pour chaque cellule :
\begin{itemize}
    \item $(H, Y{=}0)$ : $w = \frac{0{,}7 \times 0{,}7}{0{,}50} = 0{,}98$ \quad (l\'{e}g\`{e}rement diminu\'{e}, surrepr\'{e}sent\'{e})
    \item $(H, Y{=}1)$ : $w = \frac{0{,}7 \times 0{,}3}{0{,}20} = 1{,}05$ \quad (l\'{e}g\`{e}rement augment\'{e})
    \item $(F, Y{=}0)$ : $w = \frac{0{,}3 \times 0{,}7}{0{,}20} = 1{,}05$ \quad (femmes \og bonnes \fg{} : sous-repr.)
    \item $(F, Y{=}1)$ : $w = \frac{0{,}3 \times 0{,}3}{0{,}10} = 0{,}90$ \quad (femmes en d\'{e}faut : surrepr.)
\end{itemize}

\noindent Apr\`{e}s pond\'{e}ration, $S \perp Y$ : la corr\'{e}lation entre genre et d\'{e}faut dispara\^{i}t dans la perte d'entra\^{i}nement.

\begin{leconbox}[Direction g\'{e}n\'{e}rale]
Privil\'{e}gi\'{e} ``mauvais'' et non-privil\'{e}gi\'{e} ``bon'' re\c{c}oivent $w > 1$. Privil\'{e}gi\'{e} ``bon'' et non-privil\'{e}gi\'{e} ``mauvais'' re\c{c}oivent $w < 1$. C'est exactement ce qu'il faut pour casser la corr\'{e}lation $S$--$Y$.
\end{leconbox}

On entra\^{i}ne ensuite un \textbf{nouveau mod\`{e}le} avec ces poids (param\`{e}tre \texttt{sample\_weight} dans \texttt{train\_logreg}). Un mod\`{e}le repondéré est entra\^{i}n\'{e} \textbf{par attribut sensible}.

% ══════════════════════════════════════════════════════════════════════════════
\section{Seuils par groupe : correction en post-traitement}
\label{sec:postproc}

\begin{questionbox}[\'{E}tape 7 --- Peut-on corriger sans r\'{e}-entra\^{i}ner ?]
Le post-traitement ajuste le seuil de classification \textbf{s\'{e}par\'{e}ment par groupe} pour \'{e}galiser le taux de s\'{e}lection (Demographic Parity).
\end{questionbox}

\subsection{Principe}

Au lieu d'un seuil global $\tau$, on calibre $\tau_{g_0}$ et $\tau_{g_1}$ sur la validation pour que :
$$P(\hat{Y}{=}1 \mid S{=}g_0) \approx P(\hat{Y}{=}1 \mid S{=}g_1)$$

\subsection{Illustration intuitive}

\begin{center}
\begin{tikzpicture}[
    axis/.style={->, thick},
    score/.style={draw, fill=accent!10, minimum height=0.4cm, minimum width=0.15cm},
    threshold/.style={dashed, thick, alertc}
]
% Axe pour le groupe privilégié
\draw[axis] (0,2.5) -- (8,2.5) node[right, font=\small] {score};
\node[font=\small\bfseries, left] at (0,2.5) {Adulte};
\foreach \x in {0.5, 1.2, 1.5, 1.8, 2.5, 3.0, 3.5, 4.0, 4.8, 5.5, 6.2, 7.0} {
    \draw[score] (\x,2.5) -- (\x,2.9);
}
\draw[threshold] (4.2,2.0) -- (4.2,3.2) node[above, font=\scriptsize, text=alertc] {$\tau_{\text{adulte}} = 0{,}42$};

% Axe pour le groupe non-privilégié
\draw[axis] (0,0.5) -- (8,0.5) node[right, font=\small] {score};
\node[font=\small\bfseries, left] at (0,0.5) {Jeune};
\foreach \x in {1.0, 1.8, 2.5, 3.2, 4.5, 5.5, 6.5} {
    \draw[score] (\x,0.5) -- (\x,0.9);
}
\draw[threshold] (3.8,0.0) -- (3.8,1.2) node[above, font=\scriptsize, text=alertc] {$\tau_{\text{jeune}} = 0{,}38$};

% Légende
\node[font=\scriptsize, text=gray] at (4,-0.5) {Seuil plus bas pour le groupe non-privil\'{e}gi\'{e} $\Rightarrow$ taux de s\'{e}lection \'{e}galis\'{e}};
\end{tikzpicture}
\end{center}

Les deux groupes ont des distributions de scores diff\'{e}rentes. En appliquant le m\^{e}me seuil, on s\'{e}lectionnerait des proportions diff\'{e}rentes. La calibration par groupe r\'{e}\'{e}quilibre cela.

\begin{multicols}{2}
\begin{leconbox}[Avantage]
Pas de r\'{e}-entra\^{i}nement. On utilise le m\^{e}me mod\`{e}le, seul le seuil change. Tr\`{e}s rapide.
\end{leconbox}

\begin{attentionbox}[Limite]
Calibr\'{e} sur $\sim$200 exemples : seuils potentiellement instables. Et juridiquement, des seuils diff\'{e}rents par groupe peuvent \^{e}tre probl\'{e}matiques.
\end{attentionbox}
\end{multicols}

% ══════════════════════════════════════════════════════════════════════════════
\section{Comparaison : 4 configurations}
\label{sec:comparison}

\begin{questionbox}[\'{E}tape 8 --- Quel est l'effet r\'{e}el des corrections ?]
On compare syst\'{e}matiquement 4 configurations pour chaque attribut sensible : (1) Baseline, (2) Reweighing, (3) Baseline + post-processing, (4) Reweighing + post-processing.
\end{questionbox}

\subsection{Tableau de r\'{e}sultats}

\begin{table}[H]
\centering
\small
\begin{tabular}{llcccc}
\toprule
\textbf{Attr.} & \textbf{Configuration} & \textbf{AUC} & \textbf{BalAcc} & $\boldsymbol{|\Delta\textbf{DP}|}$ & $\boldsymbol{|\Delta\textbf{EO}|}$ \\
\midrule
\rowcolor{accent!5}
Genre & Baseline          & 0.785 & 0.669 & 0.010 & 0.058 \\
\rowcolor{accent!5}
Genre & Reweighing        & 0.785 & 0.665 & 0.006 & 0.050 \\
\rowcolor{accent!10}
Genre & Baseline + PP     & 0.785 & 0.669 & 0.105 & 0.196 \\
\rowcolor{accent!10}
Genre & Reweighing + PP   & 0.785 & 0.669 & 0.128 & 0.230 \\
\midrule
\rowcolor{fair!5}
\^{A}ge  & Baseline          & 0.785 & 0.669 & 0.125 & 0.066 \\
\rowcolor{fair!5}
\^{A}ge  & Reweighing        & 0.784 & 0.640 & 0.132 & 0.098 \\
\rowcolor{fair!10}
\^{A}ge  & Baseline + PP     & 0.785 & 0.644 & \textbf{0.060} & \textbf{0.135} \\
\rowcolor{fair!10}
\^{A}ge  & Reweighing + PP   & 0.784 & 0.650 & 0.069 & 0.221 \\
\bottomrule
\end{tabular}
\caption{Performance et \'{e}quit\'{e} sur le test (PP = post-processing par seuils/groupe).}
\label{tab:comparison}
\end{table}

\subsection{Analyse d\'{e}taill\'{e}e}

\paragraph{Genre --- biais quasi nul, le PP cr\'{e}e un biais artificiel.}
Apr\`{e}s retrait de \texttt{personal\_status\_sex}, $|\Delta\text{DP}| = 0.010$ : pas de discrimination directe d\'{e}tectable. Le reweighing n'a quasi aucun effet ($0.010 \to 0.006$). \textbf{En revanche}, le post-processing par seuils par groupe \textbf{cr\'{e}e} un biais artificiel ($|\Delta\text{DP}| = 0.105$, $|\Delta\text{EO}| = 0.196$). Pourquoi ? Le calibrage des seuils sur 200 exemples de validation g\'{e}n\'{e}ralise mal au test (200 exemples).

\paragraph{\^{A}ge --- biais initial significatif, compromis DP/EO clair.}
$|\Delta\text{DP}| = 0.125$ persiste apr\`{e}s retrait de \texttt{age\_in\_years}, coh\'{e}rent avec les taux de d\'{e}faut tr\`{e}s diff\'{e}rents (younger 42\,\% vs older 27\,\%). Le post-processing r\'{e}duit $|\Delta\text{DP}|$ de 0.125 \`{a} \textbf{0.060} (am\'{e}lioration), \textbf{mais} $|\Delta\text{EO}|$ \textbf{double} : $0.066 \to 0.135$. C'est l'illustration empirique parfaite du th\'{e}or\`{e}me d'impossibilit\'{e}.

\subsection{Le th\'{e}or\`{e}me d'impossibilit\'{e} en action}
\label{sec:impossibility}

\begin{attentionbox}[Th\'{e}or\`{e}me d'impossibilit\'{e} (Chouldechova 2017, Kleinberg 2016)]
Il est \textbf{math\'{e}matiquement impossible} de satisfaire simultan\'{e}ment Demographic Parity, Equal Opportunity et calibration, sauf si les taux de base sont \'{e}gaux entre groupes --- ce qui n'est pas le cas ici.
\end{attentionbox}

\paragraph{Visualisation g\'{e}om\'{e}trique du compromis.}

\begin{center}
\begin{tikzpicture}[
    axis/.style={->, thick, gray!70},
    point/.style={circle, fill=accent, draw=accent!70!black, minimum size=5pt, inner sep=0pt},
    pointfair/.style={circle, fill=fair, draw=fair!70!black, minimum size=5pt, inner sep=0pt},
    arrow/.style={-Stealth, thick, alertc}
]
% Axes
\draw[axis] (0,0) -- (6.5,0) node[right, font=\small] {$|\Delta\text{DP}|$};
\draw[axis] (0,0) -- (0,4) node[above, font=\small] {$|\Delta\text{EO}|$};
\node[font=\scriptsize, text=gray] at (-0.3, -0.3) {0};

% Frontière de Pareto
\draw[thick, accent!50, dashed] (0.3,3.5) .. controls (1.5,1.5) and (3.5,0.5) .. (5.5,0.3)
    node[right, font=\scriptsize, text=accent!70!black] {fronti\`{e}re};

% Points : configurations
\node[point] (base) at (0.6, 1.8) {};
\node[font=\scriptsize, right=1pt] at (base) {Baseline};

\node[point] (rw) at (0.6, 2.0) {};
\node[font=\scriptsize, right=1pt] at (rw) {Reweigh.};

\node[pointfair] (pp) at (0.4, 2.6) {};
\node[font=\scriptsize, above=2pt] at (pp) {RW + PP};

% Flèche illustrant le compromis
\draw[arrow, bend left=20] (rw.north) to node[right, font=\scriptsize, midway] {$\Delta\text{DP} \downarrow$ mais $\Delta\text{EO} \uparrow$} (pp.east);

% Idéal
\node[pointfair, star, fill=fair!50] (ideal) at (0,0) {};
\node[font=\scriptsize, below right=1pt and 1pt of ideal] {id\'{e}al};

% Légende des axes
\node[font=\scriptsize, text=gray] at (3, -0.4) {moins \'{e}quitable $\rightarrow$};
\node[font=\scriptsize, text=gray, rotate=90] at (-0.6, 2) {moins \'{e}quitable $\rightarrow$};
\end{tikzpicture}
\end{center}

\noindent Les configurations du mod\`{e}le se d\'{e}placent sur une \textbf{fronti\`{e}re de Pareto} : on ne peut am\'{e}liorer une dimension qu'au prix d'une autre. Le point id\'{e}al $(0, 0)$ est inaccessible.

\begin{leconbox}[\`{A} retenir]
Nos r\'{e}sultats sur l'\^{a}ge illustrent ce th\'{e}or\`{e}me : r\'{e}duire $|\Delta\text{DP}|$ \textbf{augmente} $|\Delta\text{EO}|$. Ce n'est pas un \'{e}chec --- c'est une \textbf{limite fondamentale}. En pratique, il faut choisir le crit\`{e}re selon le contexte m\'{e}tier (\og recall \fg{} pour la m\'{e}decine, \og selection rate \fg{} pour le recrutement, etc.).
\end{leconbox}

\subsection{Ce que le notebook montre}

Le graphique en \textbf{scatter plot} (AUC vs $|\Delta\text{DP}|$) visualise ce compromis : les points se d\'{e}placent vers la gauche (plus \'{e}quitable) au prix d'un l\'{e}ger recul en performance ou d'une hausse de $|\Delta\text{EO}|$.

% ══════════════════════════════════════════════════════════════════════════════
\section{Interpr\'{e}tabilit\'{e} : comprendre les d\'{e}cisions}
\label{sec:shap}

\begin{questionbox}[\'{E}tape 9 --- Pourquoi le mod\`{e}le prend-il ces d\'{e}cisions ?]
L'interpr\'{e}tabilit\'{e} est essentielle pour (1) expliquer les d\'{e}cisions aux utilisateurs, (2) d\'{e}tecter la discrimination indirecte via les proxies. Deux m\'{e}thodes compl\'{e}mentaires sont utilis\'{e}es.
\end{questionbox}

\subsection{SHAP lin\'{e}aire exact}

Pour un mod\`{e}le lin\'{e}aire, les \textbf{valeurs de Shapley} (cf.~cours interpr\'{e}tabilit\'{e}) ont une forme ferm\'{e}e exacte :

$$\phi_i(x) = w_i \cdot \left(x_i - \mathbb{E}_{\text{train}}[x_i]\right)$$

\noindent o\`{u} $w_i$ est le coefficient du mod\`{e}le et $\mathbb{E}_{\text{train}}[x_i]$ la moyenne de la feature $i$ sur le train. La contribution est sign\'{e}e : positive si la feature pousse vers le d\'{e}faut, n\'{e}gative sinon. Pour les features cat\'{e}gorielles (one-hot), on somme les contributions des dummies associ\'{e}es.

\paragraph{Top features (baseline).} Les features les plus influentes sont typiquement :
\begin{enumerate}
    \item \texttt{checking\_status} --- le statut du compte courant
    \item \texttt{duration\_in\_month} --- la dur\'{e}e du cr\'{e}dit
    \item \texttt{credit\_amount} --- le montant
    \item \texttt{credit\_history} --- l'historique
    \item \texttt{savings\_account\_bonds} --- l'\'{e}pargne
\end{enumerate}

\subsection{Importance par permutation}

M\'{e}thode compl\'{e}mentaire : on m\'{e}lange al\'{e}atoirement une colonne et on mesure la \textbf{chute d'AUC}. Si l'AUC diminue beaucoup, la feature est importante. On r\'{e}p\`{e}te 10 fois par feature pour stabiliser.

\begin{leconbox}[Convergence des deux m\'{e}thodes]
SHAP (local $\to$ global) et permutation (global) convergent sur le classement des features. C'est rassurant : les deux approches donnent la m\^{e}me image des d\'{e}cisions du mod\`{e}le.
\end{leconbox}

\subsection{D\'{e}tection de proxies (deux approches)}

Un \textbf{proxy} est une feature qui permet de pr\'{e}dire l'attribut sensible. Deux m\'{e}thodes :

\paragraph{1. Corr\'{e}lation univari\'{e}e ($|r|$).}
Pour chaque feature, on calcule la corr\'{e}lation absolue avec l'attribut sensible binaris\'{e}. Rapide mais ignore les interactions multivari\'{e}es.

\paragraph{2. R\'{e}gression logistique \texttt{features $\to$ attribut sensible} (approche Charlotte).}
On entra\^{i}ne un mod\`{e}le pour pr\'{e}dire l'attribut sensible \`{a} partir des autres features. \textbf{Si l'AUC $> 0.5$ + $\epsilon$, l'attribut est reconstructible} $\Rightarrow$ proxy multivari\'{e} actif.

\begin{table}[H]
\centering
\begin{tabular}{lcl}
\toprule
\textbf{Attribut} & \textbf{AUC proxy (multivari\'{e})} & \textbf{Interpr\'{e}tation} \\
\midrule
\textbf{Genre} & 0.69 & Signal proxy mod\'{e}r\'{e} \\
\textbf{\^{A}ge}  & \textbf{0.86} & \textbf{Tr\`{e}s fortement reconstructible} \\
\bottomrule
\end{tabular}
\caption{D\'{e}tection multivari\'{e}e de proxies (r\'{e}gression logistique).}
\end{table}

\noindent\textbf{Recommandation Charlotte :} retirer l'attribut sensible ne suffit pas. La meilleure strat\'{e}gie est de \textbf{conserver l'attribut} et de l'utiliser explicitement comme pond\'{e}ration dans une mitigation in-processing (\textit{fairness through awareness}).

\subsection{SHAP : avant vs apr\`{e}s retrait des attributs sensibles}

Pour visualiser concr\`{e}tement le m\'{e}canisme de proxy, on entra\^{i}ne un mod\`{e}le \textbf{avec} \texttt{age\_in\_years} et \texttt{personal\_status\_sex} en plus des autres features, et on compare l'importance SHAP par feature avec le mod\`{e}le \textbf{sans} (notre baseline).

\begin{table}[H]
\centering
\begin{tabular}{lc}
\toprule
\textbf{Mod\`{e}le} & \textbf{AUC test} \\
\midrule
Avec attributs sensibles & 0.781 \\
Sans attributs sensibles (baseline) & 0.785 \\
\midrule
Co\^{u}t du retrait & $-0.004$ (n\'{e}gligeable) \\
\bottomrule
\end{tabular}
\caption{Co\^{u}t en performance du retrait des attributs sensibles.}
\end{table}

\textbf{Top 5 features qui absorbent l'effet apr\`{e}s retrait} ($\Delta$ Mean |SHAP| = sans $-$ avec, plus c'est haut plus la feature \og{} h\'{e}rite \fg{} de l'effet des attributs sensibles) :

\begin{table}[H]
\centering
\begin{tabular}{lccc}
\toprule
\textbf{Feature} & \textbf{Avec} & \textbf{Sans} & \textbf{$\Delta$} \\
\midrule
\texttt{present\_employment\_since}   & 0.171 & 0.207 & \textbf{+0.036} \\
\texttt{telephone}                    & 0.216 & 0.238 & +0.022 \\
\texttt{credit\_history}              & 0.470 & 0.489 & +0.020 \\
\texttt{duration\_in\_month}          & 0.278 & 0.289 & +0.011 \\
\texttt{number\_of\_existing\_credits} & 0.140 & 0.150 & +0.009 \\
\bottomrule
\end{tabular}
\caption{Features absorbant l'effet des attributs sensibles (top 5 par $\Delta$ SHAP).}
\end{table}

\begin{leconbox}[Le\c{c}on]
SHAP rend visible ce que la r\'{e}gression proxy mesure num\'{e}riquement. Quand on retire l'attribut sensible, son influence ne dispara\^{i}t pas : elle se \textbf{redistribue} aux features les plus corr\'{e}l\'{e}es. C'est l'illustration la plus claire de la \textbf{discrimination indirecte}.
\end{leconbox}

\begin{leconbox}[Le\c{c}on fondamentale]
La d\'{e}tection de proxies \textbf{explique} pourquoi le reweighing fonctionne pour l'\^{a}ge mais pas pour le genre. L'interpr\'{e}tabilit\'{e} n'est pas un bonus --- c'est un outil essentiel pour comprendre et diagnostiquer les probl\`{e}mes d'\'{e}quit\'{e}.
\end{leconbox}

% ══════════════════════════════════════════════════════════════════════════════
\section{Quantification d'incertitude : peut-on faire confiance aux chiffres ?}
\label{sec:incertitude}

\begin{questionbox}[\'{E}tape 10 --- Les biais mesur\'{e}s sont-ils statistiquement significatifs ?]
Sur conseil de l'enseignante (Charlotte Laclau), nous rempla\c{c}ons la robustesse adversariale par une \textbf{quantification d'incertitude par bootstrap} (cf. cours \textit{Robustness and Uncertainty}, G.~Franchi). Le bootstrap mesure la stabilit\'{e} des d\'{e}cisions face \`{a} la variabilit\'{e} des donn\'{e}es.
\end{questionbox}

\subsection{Protocole bootstrap}

On entra\^{i}ne $B = 200$ mod\`{e}les sur des \'{e}chantillons tir\'{e}s avec remplacement du jeu d'entra\^{i}nement, puis on agr\`{e}ge les r\'{e}sultats. C'est l'approche d'\textbf{ensemble} la plus simple, analogue classique des \textit{Deep Ensembles} de Lakshminarayanan et al.\ (2017).

\begin{itemize}
    \item \textbf{Intervalles de confiance sur les m\'{e}triques} : quantiles 2.5\,\% et 97.5\,\% des $B$ valeurs.
    \item \textbf{Intervalles de pr\'{e}diction par individu} : quantiles des scores des $B$ mod\`{e}les.
    \item L'incertitude \textbf{\'{e}pist\'{e}mique} est captur\'{e}e par le d\'{e}saccord entre membres de l'ensemble.
\end{itemize}

\subsection{R\'{e}sultats}

\begin{table}[H]
\centering
\begin{tabular}{lcccc}
\toprule
\textbf{M\'{e}trique} & \textbf{Moyenne} & \textbf{$\sigma$} & \textbf{IC 2.5\,\%} & \textbf{IC 97.5\,\%} \\
\midrule
AUC                  & 0.764 & 0.019 & 0.723 & 0.795 \\
Balanced Accuracy    & 0.682 & 0.023 & 0.637 & 0.726 \\
$|\Delta\text{DP}|$ genre & 0.035 & 0.028 & 0.001 & 0.105 \\
$|\Delta\text{EO}|$ genre & 0.089 & 0.044 & 0.017 & 0.176 \\
$|\Delta\text{DP}|$ \^{a}ge  & 0.104 & 0.046 & 0.020 & 0.202 \\
$|\Delta\text{EO}|$ \^{a}ge  & 0.086 & 0.048 & 0.018 & 0.194 \\
\bottomrule
\end{tabular}
\caption{Intervalles de confiance bootstrap (B=200, IC 95\,\%).}
\end{table}

\subsection{Calibration et Platt scaling}

L'\textbf{ECE} mesure l'\'{e}cart entre confiance pr\'{e}dite et fr\'{e}quence r\'{e}elle :
$\text{ECE} = \sum_b \frac{|B_b|}{n} |\text{acc}(B_b) - \text{conf}(B_b)|$.

\begin{table}[H]
\centering
\begin{tabular}{lc}
\toprule
\textbf{Mod\`{e}le} & \textbf{ECE} \\
\midrule
Baseline (R\'{e}gression logistique) & 0.082 \\
Recalibr\'{e} (Platt scaling, \texttt{CalibratedClassifierCV}) & 0.080 \\
\bottomrule
\end{tabular}
\caption{Calibration : ECE avant/apr\`{e}s Platt scaling.}
\end{table}

\begin{leconbox}[Le\c{c}on]
\textbf{L'IC95\,\% de $|\Delta\text{DP}|$ \^{a}ge s'\'{e}tend de 0.02 \`{a} 0.20.} Le biais empirique de 0.125 est dans cet intervalle, mais l'incertitude est tr\`{e}s large (test de 200) $\Rightarrow$ \textbf{prudence m\'{e}thodologique}. Sur ce dataset, les \'{e}carts d'\'{e}quit\'{e} ne sont pas significatifs \`{a} un niveau pr\'{e}cis. La r\'{e}gression logistique est par ailleurs naturellement bien calibr\'{e}e (gain marginal du Platt scaling).
\end{leconbox}

% ══════════════════════════════════════════════════════════════════════════════
\section{Conclusion : ce qu'on a appris}
\label{sec:conclusion}

\begin{questionbox}[\'{E}tape 11 --- Quelles le\c{c}ons tirer ?]
Le parcours complet nous enseigne six le\c{c}ons fondamentales, qui relient les trois piliers du cours (\'{e}quit\'{e}, interpr\'{e}tabilit\'{e}, quantification d'incertitude).
\end{questionbox}

\begin{enumerate}[leftmargin=2em, itemsep=0.5em]
    \item \textbf{Les donn\'{e}es sont biais\'{e}es sur l'\^{a}ge}, pas sur le genre. Younger 42\,\% vs older 27\,\% de d\'{e}faut, alors que le genre montre $|\Delta\text{DP}|=0.010$ apr\`{e}s retrait de \texttt{personal\_status\_sex}.

    \item \textbf{Retirer l'attribut sensible ne suffit pas (approche Charlotte).} La r\'{e}gression logistique \texttt{features $\to$ \^{a}ge} a un \textbf{AUC = 0.86} : l'\^{a}ge est largement reconstructible depuis \texttt{housing}, \texttt{employment}, \texttt{job}\ldots Pour le genre, AUC = 0.69 (signal mod\'{e}r\'{e}).

    \item \textbf{Le reweighing modifie peu le mod\`{e}le linaire} sur ce dataset. Combin\'{e} au post-processing, il r\'{e}duit $|\Delta\text{DP}|$ \^{a}ge mais cr\'{e}e des effets non triviaux sur EO.

    \item \textbf{Le compromis DP/EO est in\'{e}vitable.} Le post-processing r\'{e}duit $|\Delta\text{DP}|$ \^{a}ge (0.125 $\to$ 0.060) mais \textbf{double} $|\Delta\text{EO}|$ (0.066 $\to$ 0.135). C'est le th\'{e}or\`{e}me d'impossibilit\'{e} --- pas un \'{e}chec de la m\'{e}thode.

    \item \textbf{L'interpr\'{e}tabilit\'{e} \'{e}claire l'\'{e}quit\'{e}.} SHAP + permutation + r\'{e}gression sur l'attribut sensible = d\'{e}tection multi-niveaux des proxies.

    \item \textbf{L'incertitude statistique est large.} L'IC95\,\% bootstrap de $|\Delta\text{DP}|$ \^{a}ge va de \textbf{0.02 \`{a} 0.20}. Sur un dataset de 1000 lignes, les conclusions sur l'\'{e}quit\'{e} doivent \^{e}tre nuanc\'{e}es par l'incertitude.
\end{enumerate}

\vspace{0.5em}

\begin{center}
\begin{tikzpicture}[
    node distance=1.5cm,
    pillar/.style={rectangle, draw, rounded corners=4pt, minimum width=3cm,
                   minimum height=1.2cm, font=\bfseries, text centered, align=center},
    link/.style={<->, thick, gray!60}
]
\node[pillar, fill=fair!15] (eq) {\'{E}quit\'{e}};
\node[pillar, fill=neutral!15, right=2.5cm of eq] (interp) {Interpr\'{e}tabilit\'{e}};
\node[pillar, fill=accent!15, below right=1.2cm and 1.25cm of eq] (rob) {Incertitude};

\draw[link] (eq) -- node[above, font=\scriptsize, text=gray] {proxies} (interp);
\draw[link] (eq) -- node[left, font=\scriptsize, text=gray, pos=0.4] {IC bootstrap} (rob);
\draw[link] (interp) -- node[right, font=\scriptsize, text=gray, pos=0.4] {confiance} (rob);
\end{tikzpicture}
\end{center}

\vspace{0.3em}
\begin{center}
\small\textit{Les trois piliers sont interd\'{e}pendants : l'interpr\'{e}tabilit\'{e} r\'{e}v\`{e}le les probl\`{e}mes d'\'{e}quit\'{e},\\la quantification d'incertitude pr\'{e}cise la confiance qu'on peut leur accorder.}
\end{center}

\subsection*{Arbre de d\'{e}cision : que faire selon la situation ?}

\begin{center}
\begin{tikzpicture}[
    node distance=0.5cm and 0.8cm,
    decision/.style={diamond, draw, fill=accent!10, aspect=2,
                     text width=3.2cm, align=center, font=\scriptsize, inner sep=2pt},
    action/.style={rectangle, draw, fill=fair!10, rounded corners=3pt,
                   text width=3.5cm, align=center, font=\scriptsize, inner sep=4pt},
    arrow/.style={-Stealth, thick}
]
\node[decision] (start) {Y a-t-il des biais\\dans les donn\'{e}es ?};
\node[action, below=1cm of start] (step1) {Mesurer les taux\\de base par groupe};
\node[decision, below=of step1] (proxy) {Y a-t-il des proxies\\corr\'{e}l\'{e}s \`{a} $S$ ?};
\node[action, below left=0.7cm and 0.5cm of proxy] (rmcol) {Retirer la colonne\\sensible suffit};
\node[action, below right=0.7cm and 0.5cm of proxy] (rw) {Reweighing\\+ post-processing};
\node[decision, below=2.3cm of proxy] (crit) {Quel crit\`{e}re\\prioriser ?};
\node[action, below left=0.7cm and 0.3cm of crit] (dp) {DP : seuils\\par groupe};
\node[action, below right=0.7cm and 0.3cm of crit] (eo) {EO : reweighing\\(\textbf{pas} PP)};

\draw[arrow] (start) -- (step1);
\draw[arrow] (step1) -- (proxy);
\draw[arrow] (proxy) -- node[above left, font=\tiny] {non} (rmcol);
\draw[arrow] (proxy) -- node[above right, font=\tiny] {oui} (rw);
\draw[arrow] (rw) |- (crit);
\draw[arrow] (crit) -- node[above left, font=\tiny] {parit\'{e}} (dp);
\draw[arrow] (crit) -- node[above right, font=\tiny] {recall} (eo);
\end{tikzpicture}
\end{center}

\vspace{1em}
\hrule
\vspace{0.5em}
\small\noindent\textit{Code : \texttt{scikit-learn} + \texttt{fairlearn} + \texttt{shap}, reproductibilit\'{e} via \texttt{papermill}.}

\end{document}
